\chapter{Preliminaries}\label{chap:preliminaries}

This chapter fixes notation and describes how we define and prove security.
It then introduces the primitives that later chapters build on: commitment
schemes, pseudorandom functions, hash functions, signature schemes, public key
encryption schemes and proof systems. Their definitions are standard and follow
Katz and Lindell~\cite{KatLin14}. Finally, it gives bilinear groups, the
assumptions we rely on, and the instantiations used by more than one chapter:
Sigma protocols and the Fiat-Shamir transform, the Pedersen commitment scheme,
the ElGamal encryption scheme, SPS schemes with $\AGHO$, and the BBS+ signature
scheme.

The thesis contains numerous primitives, protocols, security definitions,
instantiations and assumptions. To ease understanding of them, and to separate
our contributions from the prior work we build on, we use a number of different
coloured boxed environments. Abstract primitives are given directly in prose, and ideal functionalities are
displayed as plain figures. Security experiments are in green boxes, assumptions
in pink boxes, instantiations of primitives in blue boxes, and the protocols that
are contributions of this work in purple boxes. Every box and figure is numbered and titled. Once
experiments, assumptions, primitives and protocols have been introduced, later
chapters refer back to them rather than restating them, and each chapter is
otherwise as self-contained as possible.

\section{Notation}\label{sec:prelims-notation}

We write $\{0,1\}^*$ for the set of bit strings of any length, $|x|$ for the
length of $x$, and $x \| y$ for the concatenation of $x$ and $y$. Both $[n]$ and
$[1,n]$ denote the set $\{1, \ldots, n\}$. Vectors carry an arrow, as in
$\vec{x}$, and matrices are upright capitals. We write $x \sample S$ for
sampling $x$ uniformly from a set $S$, $y \sample \algo{A}(x)$ for running a
randomised algorithm $\algo{A}$ on input $x$, and $y \leftarrow \algo{A}(x)$ for
running a deterministic one. Where a security argument needs to name the
randomness, we write $\algo{A}(x; \rho)$, so that $y \sample \algo{A}(x)$ is
$y \leftarrow \algo{A}(x; \rho)$ for uniform $\rho$. A list of the notation used
throughout is given at the front of the thesis.

\section{Defining secure protocols}\label{sec:prelims-security}

The contributions of this thesis include both constructions of protocols and
constructions of their underlying primitives. We take different approaches to
defining and proving the security of each of these, described below after the
notions of efficiency and negligible probability that both rely on.

The definitions of the security goals that we wish to achieve are not complete
without defining the capabilities of the adversary and its success probability.
There are two main ways of doing this: the concrete approach and the asymptotic
approach. The concrete approach assigns numbers to the amount of operations that
the adversary is able to execute and a fixed bound to the probability of
success. The asymptotic approach instead treats security as a function
parameterised by a security parameter. Then as this parameter grows, the
resources needed to break a scheme grow faster than any fixed polynomial, and
the probability of success of the adversary becomes vanishingly small.

We have chosen to follow the asymptotic approach. This deserves some
consideration, as we are explicitly aiming for concrete efficiency in the design
of our protocols and primitives. We wish for the schemes to be used in practice.
The primitives and protocols we design have elements chosen in their
instantiations specifically for the efficiency benefits. However, we wish to
separate the instantiation choices from the design and security proofs of the
abstract constructions that we present.

An adversary is considered efficient if it runs in
probabilistic polynomial time (PPT) in the security parameter. We call this
parameter $\secparam$, and so the adversary's running time is bounded by
$p(\secparam)$ for some fixed polynomial $p$. It may use randomness in its
computation. Honest parties must also satisfy this definition of efficiency.
Every algorithm takes $\secparam$ as an input, explicitly or implicitly.

For the success probability of the adversary, we define small with a negligible
function. A function $f$ from the natural
numbers to non-negative real numbers is negligible if it is eventually below every
inverse polynomial. Concretely, for every
polynomial $p$, there exists $N$ such that $f(\secparam) < 1/p(\secparam)$ for
all $\secparam > N$. We write $\negl(\secparam)$ for an unspecified negligible
function. Negligible functions have a property of closure which carries well to
their use in security definitions in cryptography. This is that a polynomial
number of events each occurring with negligible probability, still occurs with
negligible probability in total. This means an adversary repeating an attack
with negligible success probability polynomially many times will still have a
negligible probability of success.

Indistinguishability combines these two notions. A distinguisher is a PPT
adversary that outputs a single bit. Two families of distributions indexed by
$\secparam$ are computationally indistinguishable if every distinguisher outputs
$1$ with probabilities that differ only negligibly between the two. This notion
is used later by the DDH assumption, by zero knowledge, and to compare real and
ideal executions.

By primitives, we refer to building blocks such as commitments and signature
schemes. Security in this case is handled through defining experiments, a
limited number of oracles which an adversary has access to, and then showing
that the adversary cannot win with greater probability than some predefined
amount.

An experiment is a game between a challenger, which runs the honest
algorithms, and the adversary $\adv$. The adversary may be given oracle access,
in which it submits inputs of its choice and the challenger answers them using
secrets that $\adv$ does not see. The experiment states when $\adv$ wins. The
advantage $\Adv$ of $\adv$ is its probability of winning, or, where guessing
already wins half the time, the distance of that probability from $1/2$. A
property holds if every PPT adversary has negligible advantage. Proofs are
reductions: from an adversary with non-negligible advantage, we build an
algorithm that breaks an assumption or another primitive. Many proofs go
through a sequence of games, and bound the change in the success probability of
$\adv$ between consecutive games.

For larger protocols with many parties and interfaces, we capture all security
goals at once inside an ideal functionality. The ideal functionality represents
an incorruptible trusted party that receives the inputs of each protocol task
and returns the prescribed outputs. We then construct a simulator to show that
the view of an adversary in the real world, interacting with our protocol, is
indistinguishable from the view when interacting with the ideal functionality
itself.

We have the outputs of the honest parties and~$\adv$ during an execution of a
protocol~$\Pi$ in the real world captured in
$\textsc{Exec}_{\Pi,\adv}(\secparam)$, and we have the outputs of the honest
parties and~$\simu$ interacting with an ideal functionality~$\idfu{}$ in the
ideal world captured by $\textsc{Exec}_{\idfu{},\simu}(\secparam)$. We then say
that the protocol~$\Pi$ securely realises the ideal functionality~$\idfu{}$ if
for every PPT adversary~$\adv$, there exists a PPT simulator~$\simu$ where
$\textsc{Exec}_{\Pi,\adv}(\secparam)$ is computationally indistinguishable from
$\textsc{Exec}_{\idfu{},\simu}(\secparam)$. As the ideal functionality achieves
all security properties by design, this then shows that the protocol achieves
them too.

The subscript of $\idfu{}$ names the task it captures. Each party receives its
own output from $\idfu{}$ and nothing more. The view of a party consists of its
inputs, its randomness and the messages it receives. In the ideal execution,
honest parties communicate only with $\idfu{}$, and $\simu$ takes the place of
$\adv$.

\begin{outline}
  \item we work in the standalone model, executions composed sequentially
  \item UC gives security under arbitrary concurrent composition: every
environment supplies inputs and talks to the adversary throughout, so the
simulator cannot rewind (cite FOCS:Canetti01)
  \item why it matters: several of our simulators extract witnesses by
rewinding (forking lemma), sound standalone but not in UC (your old ch05
52-definitions ``Security model'' paragraph says this)
  \item straight-line extraction would recover UC: prover encrypts its witness
under a key in the CRS and proves it, or Fischlin instead of Fiat-Shamir (cite
C:Fischlin05)
  \item both cost proof size; we only pay it for the withdrawal proof of
\Cref{chap:offline-payments}
\end{outline}

In both kinds of definition, the adversary corrupts some participants.
Honest-but-curious, or semi-honest, participants follow the protocol, and the
adversary sees their views and combines them with the views of the other
corrupted participants. Malicious participants deviate from the protocol
arbitrarily. Under static corruption, the set of corrupted participants is fixed
before the protocol starts. Under adaptive corruption, the adversary may corrupt
further participants during the execution, based on what it has seen. Each
chapter states which participants may be corrupted, and how.

\section{Commitment schemes}\label{sec:prelims-commitments}

Commitment schemes let parties commit to messages while keeping them hidden until
opened.

\paragraph{Syntax.} A commitment scheme consists of algorithms defined as
follows:
%
\begin{description}

  \item[$\ck \sample \ComGen(\secparam)$ :] The key generation algorithm
takes the security parameter as input, and outputs a commitment key $\ck$.

  \item[$\com \leftarrow \ComCommit(\ck, \msg, \openinfo)$ :] The commitment
algorithm outputs commitment $\com$ to message $\msg \in \msgspace$ using key $\ck$ and
opening information $\openinfo$, which the committer samples uniformly and keeps
until the commitment is opened.

\end{description}
%
A committer opens $\com$ by revealing $\msg$ and $\openinfo$, and the receiver
checks that $\com = \ComCommit(\ck, \msg, \openinfo)$. Commitment schemes
satisfy hiding and binding, which we define below.

\paragraph{Hiding.} Hiding requires that a commitment reveal no information
about its message before it is opened, which \Cref{exp:hiding} captures.

\begin{experimentbox}{Commitment schemes: hiding}{hiding}

  \begin{enumerate}

    \item The challenger generates $\ck \sample \ComGen(\secparam)$ and sends
$\ck$ to $\adv$.

    \item The adversary outputs two messages $(\msg_0, \msg_1)$.

    \item The challenger samples $b \sample \{0,1\}$ and fresh opening
information $\openinfo$, computes the challenge commitment $\com^\star
\leftarrow \ComCommit(\ck, \msg_b, \openinfo)$, and sends
$\com^\star$ to $\adv$.

    \item The adversary outputs a bit $b'$.

  \end{enumerate}

  $\adv$ wins the experiment if $b' = b$.

\end{experimentbox}

The hiding advantage of $\adv$ is defined as follows:
%
\begin{equation*}
%
  \Adv_{\Com,\adv}^{\mathsf{hiding}}(\secparam) = \left| \Pr[\adv\text{
wins}]-\frac{1}{2} \right|.
%
\end{equation*}
%
We say the commitment scheme $\Com$ is hiding if this advantage is negligible for
every PPT adversary $\adv$.

\paragraph{Binding.} We also require commitment schemes to satisfy the binding
property of \Cref{exp:binding}, which prevents an adversary from opening the same
commitment to two different messages.

\begin{experimentbox}{Commitment schemes: binding}{binding}

  \begin{enumerate}

    \item The challenger generates $\ck \sample \ComGen(\secparam)$ and sends
$\ck$ to $\adv$.

    \item The adversary outputs $(\com, \msg, \openinfo, \msg', \openinfo')$.

  \end{enumerate}

  $\adv$ wins the experiment if the following hold:
  %
  \begin{equation*}
  %
    \msg \neq \msg', \quad \ComCommit(\ck, \msg, \openinfo) = \com =
\ComCommit(\ck, \msg', \openinfo').
  %
  \end{equation*}

\end{experimentbox}

The binding advantage of $\adv$ is defined as follows:
%
\begin{equation*}
%
  \Adv^{\mathsf{binding}}_{\Com,\adv}(\secparam) = \Pr[\adv\text{ wins}].
%
\end{equation*}
%
A commitment scheme is binding if this advantage is negligible for every PPT
adversary $\adv$.

\section{Pseudorandom functions}\label{sec:prelims-prfs}

Pseudorandom functions enable a party holding a secret key to evaluate a
function whose outputs are indistinguishable from those of a truly random
function to anyone without the key.

\paragraph{Syntax.} A pseudorandom function family consists of algorithms
defined as follows:
%
\begin{description}

  \item[$(\pp, \key) \sample \PRFgen(\secparam)$ :] The key generation algorithm
takes the security parameter as input, and outputs public parameters $\pp$
together with a key $\key \in \keyspace$.

  \item[$y \leftarrow \PRFeval(\pp, \key, x)$ :] The evaluation algorithm takes
public parameters $\pp$, key $\key$, and input $x \in \domain$, and outputs
$y \in \range$.

\end{description}
%
Here $\domain$ and $\range$ are the domain and range of the family. We say
that the pair of algorithms $\PRF = (\PRFgen, \PRFeval)$ is a pseudorandom
function family if $\PRFeval(\pp, \key, \cdot)$ is
efficiently computable for every $(\pp, \key)$ output by $\PRFgen(\secparam)$,
and if it satisfies pseudorandomness, which we define below. We use the standard
syntax and security notions for pseudorandom functions~\cite{KatLin14}.

The public parameters are carried explicitly because the algebraic
instantiations of later chapters need them: the domain, the range, or the
description of the function itself may depend on values that are common to every
key, and those values are public. The textbook setting is the special
case of empty $\pp$, and we suppress $\pp$ where it plays no role.

\paragraph{Pseudorandomness.} Pseudorandomness requires that no efficient
adversary, given the public parameters and oracle access to the function, can
distinguish it from a function sampled uniformly at random from all functions
$\domain \to \range$. \Cref{exp:prf} gives the experiment.

\begin{experimentbox}{Pseudorandom functions: pseudorandomness}{prf}

  \begin{enumerate}

    \item The challenger samples $b \sample \{0,1\}$ and generates $(\pp, \key)
\sample \PRFgen(\secparam)$. If $b = 1$, it additionally samples a function $R$
uniformly at random from all functions $\domain \to \range$. It sends $\pp$ to
the adversary.

    \item The adversary is given oracle access to $\mathcal{O}_{\PRF}$, which on
input $x \in \domain$ returns $\PRFeval(\pp, \key, x)$ if $b = 0$, and $R(x)$ if
$b = 1$.

    \item The adversary outputs a bit $b'$.

  \end{enumerate}

  $\adv$ wins the experiment if $b' = b$.

\end{experimentbox}

The pseudorandomness advantage of $\adv$ is defined as follows:
%
\begin{equation*}
%
  \Adv_{\PRF,\adv}^{\mathsf{prf}}(\secparam) = \left| \Pr[\adv\text{
wins}]-\frac{1}{2} \right|.
%
\end{equation*}
%
We say $\PRF$ is a secure pseudorandom function if this advantage is negligible for
every PPT adversary $\adv$ making a polynomial number of queries to
$\mathcal{O}_{\PRF}$.

Pseudorandomness also gives a form of hiding. To anyone without the key, outputs
on distinct inputs are indistinguishable from independent random values, and so
reveal nothing about the key or about each other.

\section{Hash functions}\label{sec:prelims-hash}

Hash functions compress inputs of arbitrary length into a short digest. They
are usually defined as keyed families, with collision resistance, preimage
resistance and second preimage resistance as the standard security
notions~\cite{KatLin14}. Our constructions do not rely on these properties
directly. Instead, our security proofs model hash functions as random oracles, and the
one hash function used inside a circuit, MiMC, is assumed to be a pseudorandom
function (\Cref{sec:txid-hash-instantiation}). We write $\hash(x)$ for the
hash of $x$.

The range can be any set, and later sections also hash
into fields and groups. Schemes that use several hash functions write them
$\hashone, \hashtwo, \ldots$ and give each its own domain separation tag. A
domain separation tag is a fixed string prefixed to every input, and distinct
tags make the hash functions behave as independent functions.

\paragraph{One-way functions.} A function $f$ is one-way if it is efficiently
computable and, for a uniform input $x$, no PPT adversary given $f(x)$ outputs
any $x'$ with $f(x') = f(x)$ with non-negligible probability.

Our instantiations use exponentiation $x \mapsto \gen^x$ in a group of prime order
as the one-way function, which is one-way under the DL assumption in
\Cref{asm:dl}.

\paragraph{The random oracle model.} Our security
proofs generally model a hash
function as a random oracle: a function sampled uniformly at random from all
functions with the given domain and range, to which every party, including the
adversary, has oracle access. The random oracle answers each fresh query with an
independently and uniformly sampled output from the function's range, records
the resulting input and output pair, and returns the same output on repeated
queries. In security proofs, a reduction or simulator answers the adversary's
oracle queries itself, so it sees every query, and may choose the output
assigned to particular queries. We refer to this as programming the random
oracle. Security proven this way holds in the random oracle model (ROM).

\section{Digital signature schemes}\label{sec:prelims-signatures}

Digital signature schemes enable signers to produce publicly verifiable
signatures on messages.

\paragraph{Syntax.} A digital signature scheme consists of the following
algorithms:

\begin{description}

  \item[$(\sk,\pk) \sample \DSgen(\secparam)$ : ] The key generation algorithm
outputs secret signing key $\sk$ and public verification key $\pk$.

  \item[$\signature \sample \DSsign(\sk, \msg)$ : ] The signing algorithm signs
message $\msg \in \msgspace$ using secret key $\sk$, giving output $\signature$.

  \item[$\{0,1\} \leftarrow \DSverify(\pk,\msg,\signature)$ : ] The verification
algorithm verifies signature $\signature$ on message $\msg$ using public key
$\pk$, giving output $0$ or $1$.

\end{description}

Digital signature schemes satisfy correctness and unforgeability, one form of
which we define below. Again, we use standard syntax and security notions for
digital signature schemes~\cite{KatLin14}.

\paragraph{Correctness.} A digital signature scheme is correct if, for every key
pair $(\sk,\pk)$ generated by $\DSgen(\secparam)$ and every message $\msg$, the
following holds:
%
\begin{equation*}
%
  \DSverify(\pk, \msg, \DSsign(\sk, \msg)) = 1.
%
\end{equation*}
%

\paragraph{Unforgeability.} We define existential unforgeability under chosen
message attack (EUF-CMA) using \Cref{exp:euf-cma}.

\begin{experimentbox}{Digital signature schemes: EUF-CMA}{euf-cma}

  \begin{enumerate}

    \item The challenger generates $(\sk,\pk) \sample \DSgen(\secparam)$ and
sends $\pk$ to $\adv$.

    \item The adversary is given oracle access to $\DSsigno$, which on input a
message $\msg$ returns $\signature = \DSsign(\sk, \msg)$ and adds $\msg$ to
the set $\mssg$, which is initially empty.

    \item The adversary outputs $(\msg^\star, \signature^\star)$.

  \end{enumerate}

  $\adv$ wins the experiment if the following hold:
  %
  \begin{equation*}
  %
    \DSverify(\pk,\msg^\star,\signature^\star) = 1 \ \wedge\ \msg^\star \notin
\mssg.
  %
  \end{equation*}

\end{experimentbox}

The advantage of $\adv$ is defined as follows:
%
\begin{equation*}
%
  \Adv^{\mathsf{euf\text{-}cma}}_{\DS, \adv}(\secparam) = \Pr[\adv\text{
wins}].
%
\end{equation*}
%
We say the digital signature scheme $\DS$ is EUF-CMA secure if this advantage is
negligible for every PPT adversary $\adv$.

Strong unforgeability under chosen message attack (sUF-CMA) is defined by the
same experiment except that the oracle stores the pairs $(\msg, \signature)$ it
returns, and any valid pair that was not returned counts as a forgery. Every
scheme in this thesis needs only EUF-CMA security of the signature schemes it
builds on.

\section{Public key encryption schemes}\label{sec:prelims-encryption}

Public key encryption schemes enable confidential communication by allowing
anyone possessing a public key to encrypt messages that can only be decrypted
using the corresponding secret key.

\paragraph{Syntax.} A public key encryption scheme consists of the following
algorithms:

\begin{description}

  \item[$(\sk, \pk) \sample \PKEgen(\secparam)$ : ] The key generation algorithm
outputs secret decryption key $\sk$ and public encryption key $\pk$.

  \item[$\ciph \sample \PKEencrypt(\pk, \msg)$ : ] The encryption algorithm
encrypts message $\msg \in \msgspace$ using public key $\pk$, giving output
ciphertext
$\ciph$.

  \item[$\msg \leftarrow \PKEdecrypt(\sk, \ciph)$ : ] The decryption algorithm
decrypts ciphertext $\ciph$ using secret key $\sk$, giving output message
$\msg$.

\end{description}

Public key encryption schemes satisfy correctness and indistinguishability. We
define correctness and indistinguishability under chosen-plaintext attack
below. We use the standard syntax and security notions for public key
encryption~\cite{KatLin14}.

\paragraph{Correctness.} A public key encryption scheme is correct if, for every
key pair $(\sk,\pk)$ generated by $\PKEgen(\secparam)$ and every message $\msg$,
the following holds:
%
\begin{equation*}
%
  \PKEdecrypt(\sk, \PKEencrypt(\pk, \msg)) = \msg.
%
\end{equation*}
%
\paragraph{IND-CPA security.} We define indistinguishability under
chosen-plaintext attack (IND-CPA) security using \Cref{exp:ind-cpa}.

\begin{experimentbox}{Public key encryption schemes: IND-CPA}{ind-cpa}

  \begin{enumerate}

    \item The challenger generates $(\sk, \pk) \sample \PKEgen(\secparam)$ and
sends $\pk$ to $\adv$.

    \item The adversary outputs two equal length messages $(\msg_0,\msg_1)$.

    \item The challenger samples $b \sample \{0,1\}$, computes the challenge
ciphertext $\ciph^\star \sample \PKEencrypt(\pk,\msg_b)$, and sends
$\ciph^\star$ to the adversary.

    \item The adversary outputs a bit $b'$.

  \end{enumerate}

  $\adv$ wins the experiment if $b' = b$.

\end{experimentbox}

The IND-CPA advantage of $\adv$ is defined as follows:
%
\begin{equation*}
%
  \Adv^{\mathsf{ind\text{-}cpa}}_{\PKE, \adv}(\secparam) = \left|
\Pr[\adv\text{ wins}] - \frac{1}{2} \right|.
%
\end{equation*}
%
We say the public key encryption scheme $\PKE$ is IND-CPA secure if this
advantage is negligible for every PPT adversary $\adv$.

Indistinguishability under chosen-ciphertext attack (IND-CCA) is defined by the
same experiment, with the adversary also given an oracle that decrypts any
ciphertext other than $\ciph^\star$. IND-CCA security implies IND-CPA security.

\section{Proof systems}\label{sec:prelims-nizks}

Zero-knowledge proof systems enable provers to convince verifiers that a
statement is true without revealing any information beyond the validity of the
statement. Soundness requires that verifiers are not convinced of false
statements. We use non-interactive proof systems throughout~\cite{STOC:BluFelMic88}.

A relation $\rel$ is an efficiently decidable set of pairs $(\inst, \witness)$
of a statement and a witness. We write $\rel(\inst, \witness) = 1$ when
$(\inst, \witness)$ is in $\rel$.

\paragraph{Syntax.} We call a non-interactive proof system a non-interactive
zero-knowledge (NIZK) proof system when it is also zero knowledge, as defined
below. We write $\NIZK$ throughout, as all of our constructions satisfy zero
knowledge. A non-interactive proof system $\NIZK$ consists of the following
algorithms.

\begin{description}

  \item[$\crs \sample \NIZKgen(\secparam, \rel)$ :] The setup algorithm, on
input the security parameter and a relation $\rel$, outputs a common
reference string $\crs$, public parameters that prover and verifier both use.

  \item[$\algoproof \sample \NIZKprove(\crs, \inst, \witness)$ :] The proving
algorithm takes as input a common reference string $\crs$, statement $\inst$,
and witness $\witness$, and outputs a proof $\algoproof$.

  \item[$\{0, 1\} \leftarrow \NIZKverify(\crs, \inst, \algoproof)$ :] The
verification algorithm takes as input a common reference string $\crs$,
statement $\inst$, and proof $\algoproof$, and outputs $0$ or $1$.

\end{description}

\paragraph{Correctness.} A non-interactive proof system must satisfy the
following correctness requirement: for every common reference string
$\crs$ generated by $\NIZKgen(\secparam, \rel)$ and every statement and witness
pair
$(\inst, \witness)$ with $\rel(\inst, \witness) = 1$, the following holds:
%
\begin{equation*}
%
  \NIZKverify(\crs, \inst, \NIZKprove(\crs, \inst, \witness)) = 1.
%
\end{equation*}
%
\paragraph{Soundness.} A non-interactive proof system is sound if, for every
PPT adversary $\adv$, the following probability is negligible:
%
\begin{equation*}
%
  \Pr\left[ \begin{array}{l} \NIZKverify(\crs, \inst, \algoproof) = 1 \\
\wedge\ \nexists \witness : \rel(\inst, \witness) = 1 \end{array} \
\middle| \ \begin{array}{l} \crs \sample \NIZKgen(\secparam, \rel) \\
(\inst, \algoproof) \sample \adv(\crs) \end{array} \right].
%
\end{equation*}
%
The constructions in this thesis need stronger notions than soundness. We define
zero knowledge and simulation extractability below, and introduce any
construction-specific requirements in the chapters that use them.

\paragraph{Labelled proofs.} Certain constructions require proofs that are bound
to a message to prevent proof replaying.  We can attach a label to each proof to
prevent this.  A label is an arbitrary string that is fixed at proving time.  A
labelled proof system extends the syntax as follows:

\begin{description}

  \item[$\algoproof \sample \NIZKprove(\crs, \inst, \witness, \lbl)$ :] The
proving algorithm also takes a label $\lbl \in \{0,1\}^*$.

  \item[$\{0, 1\} \leftarrow \NIZKverify(\crs, \inst, \algoproof, \lbl)$ :] The
verification algorithm also takes the label $\lbl$, and accepts only
proofs produced under that label.

\end{description}

Correctness is then defined as above, over all labels. A labelled proof system that is
zero knowledge and simulation extractable, as defined below, is what Chase and
Lysyanskaya call a signature of knowledge, with the label as the signed
message~\cite{C:ChaLys06}.

Zero knowledge and simulation extractability require defining further
algorithms: simulators, producing proofs without witnesses, and extractors,
recovering witnesses from proofs. We define the algorithms as follows:

\begin{description}

  \item[$(\crs, \td) \sample \NIZKsimgen(\secparam, \rel)$ :] The simulated setup
algorithm outputs a common reference string $\crs$ together with a trapdoor
$\td$.  The distribution of $\crs$ is computationally indistinguishable from
that produced by $\NIZKgen(\secparam, \rel)$.

  \item[$\algoproof \sample \NIZKsimprove(\crs, \td, \inst, \lbl)$ :] The
simulator takes the trapdoor, a statement $\inst$, and a label $\lbl$, and
outputs a proof, without access to a witness.

  \item[$\witness \leftarrow \NIZKextract(\crs, \td, \inst, \algoproof, \lbl)$
:] The extractor takes the trapdoor, a statement, a proof, and a label, and
outputs a candidate witness.

\end{description}

\paragraph{Zero knowledge.} A labelled proof system is zero knowledge if simulated
proofs are indistinguishable from honestly generated ones, as defined in
\Cref{exp:zk}.

\begin{experimentbox}{Proof systems: zero knowledge}{zk}

  \begin{enumerate}

    \item The challenger samples $b \sample \{0,1\}$.  If $b = 0$ it generates
$\crs \sample \NIZKgen(\secparam, \rel)$. If $b = 1$ it generates $(\crs, \td)
\sample \NIZKsimgen(\secparam, \rel)$.  It sends $\crs$ to $\adv$.

    \item $\adv$ is given oracle access to $\simoracle$, which on input $(\inst,
\witness, \lbl)$ with $\rel(\inst, \witness) = 1$ returns $\NIZKprove(\crs,
\inst, \witness, \lbl)$ if $b = 0$, and $\NIZKsimprove(\crs, \td, \inst, \lbl)$
if $b = 1$.

    \item $\adv$ outputs a bit $b'$.

  \end{enumerate}

  $\adv$ wins the experiment if $b' = b$.

\end{experimentbox}

The advantage of $\adv$ is defined as follows:
%
\begin{equation*}
%
  \Adv^{\mathsf{zk}}_{\NIZK, \adv}(\secparam) = \left| \Pr[\adv\text{
wins}] - \frac{1}{2} \right|.
%
\end{equation*}
%
We say the proof system satisfies the property of zero knowledge if this
advantage is negligible for every PPT adversary $\adv$.

\paragraph{Simulation extractability.} Our security arguments need a
witness to be recoverable from all accepting proofs, even after an adversary has
observed simulated proofs. Simulation extractability, defined by
\Cref{exp:simext}, captures this~\cite{FOCS:Sahai99,C:DDOPS01,AC:Groth06}.

\begin{experimentbox}{Proof systems: simulation extractability}{simext}

  \begin{enumerate}

    \item The challenger generates $(\crs, \td) \sample \NIZKsimgen(\secparam,
\rel)$
and sends $\crs$ to $\adv$.

    \item $\adv$ is given oracle access to $\simoracle$, which on input $(\inst,
\lbl)$ returns $\algoproof \sample \NIZKsimprove(\crs, \td, \inst, \lbl)$ and
stores $(\inst, \lbl)$ in the set $\Query$.

    \item $\adv$ outputs $(\inst^*, \algoproof^*, \lbl^*)$.

    \item Challenger extracts $\witness^* \leftarrow \NIZKextract(\crs, \td,
\inst^*, \algoproof^*, \lbl^*)$.

  \end{enumerate}

  $\adv$ wins the experiment if the following hold:
  %
  \begin{equation*}
  %
    \NIZKverify(\crs, \inst^*, \algoproof^*, \lbl^*) = 1 \ \wedge\ (\inst^*,
\lbl^*) \notin \Query \ \wedge\ \rel(\inst^*, \witness^*) \neq 1.
  %
  \end{equation*}

\end{experimentbox}

The simulation extractability advantage of $\adv$ is defined as follows:
%
\begin{equation*}
%
  \Adv^{\mathsf{simext}}_{\NIZK, \adv}(\secparam) = \Pr[\adv\text{
wins}].
%
\end{equation*}
%
A labelled proof system is simulation extractable if this advantage is negligible
for every PPT adversary $\adv$.

Simulation extractability is a strong form of
non-malleability, which prevents an adversary from turning proofs it has seen
into proofs of related statements~\cite{AC:Groth06}. The condition $(\inst^*, \lbl^*) \notin \Query$ still allows $\adv$ to modify a
simulated proof into another proof for the same statement and label. As a
witness can be extracted, simulation extractable proof systems are proofs of
knowledge.

\section{Instantiations}\label{sec:prelims-inst}

\begin{outline}
  \item what this section gives: the groups, assumptions and concrete schemes the rest builds on
  \item what is included here: the standard instantiation of a primitive from this chapter (Pedersen, ElGamal), or prior work used in more than one chapter ($\AGHO$, BBS+)
  \item schemes used in one chapter only, variants modified for one chapter, and our own constructions are in the chapter that uses them
\end{outline}

\subsection{Groups and assumptions}\label{sec:prelims-groups}

\begin{outline}
  \item $\Zp$ for prime $p$: integers mod $p$, elements called field elements; $\Zp^* = \Zp \setminus \{0\}$
  \item cyclic groups of prime order $p$, written multiplicatively; their elements are group elements
  \item discrete logarithm: for generator $\gen$ and $a = \gen^x$, $x \in \Zp$ is the discrete logarithm of $a$ relative to $\gen$
\end{outline}

Most of our schemes use three such groups, $\group{1}$, $\group{2}$, and
$\group{T}$ of prime order $p$, with generators $\gen \in \group{1}$ and $\ggen
\in \group{2}$. The groups are equipped with an efficiently computable pairing
$\pairing \colon \group{1} \times \group{2} \rightarrow \group{T}$. For all $a
\in \group{1}$, $\widetilde{b} \in \group{2}$, and $x,y \in \Zp$, bilinearity
means that the following holds:
%
\begin{equation*}
%
  \pairing(a^x,\widetilde{b}^y) = \pairing(a,\widetilde{b})^{xy}.
%
\end{equation*}
%
The pairing is non-degenerate, so $\pairing(\gen,\ggen)$ generates $\group{T}$.
We call $\group{1}$ and $\group{2}$ the source groups and $\group{T}$ the target
group. We work in the asymmetric, or type III, setting, in which $\group{1} \neq
\group{2}$ and no efficiently computable homomorphism between them is known.
Elements of $\group{2}$ carry a tilde, and field elements are usually written
with Greek letters.

Bilinearity allows relations between exponents in the source groups to be
checked through equalities in $\group{T}$. The equality
$\pairing(a,\ggen) = \pairing(\gen,\widetilde{b})$ shows that $a$ and
$\widetilde{b}$ have the same discrete logarithm relative to their generators,
without that discrete logarithm being recovered. A pairing product equation is
an equation of the form $\prod_j \pairing(a_j, \widetilde{b}_j) = t$ for $t \in
\group{T}$. Our constructions use such equations to verify signatures and
algebraic relations across $\group{1}$ and $\group{2}$.

\begin{outline}
  \item we rely on hardness of discrete logarithms (\Cref{asm:dl}) and on DDH (\Cref{asm:ddh}), cite DifHel76, Boneh98
  \item assumptions specific to one chapter are given where they are used
\end{outline}

\begin{assumptionbox}{Discrete logarithm (DL)}{dl}

  The DL assumption holds in a group of prime order $p$ with generator $\gen$
if, for $x \sample \Zp$, no PPT adversary given $\gen^x$ outputs $x$ with
non-negligible probability.

\end{assumptionbox}

\begin{assumptionbox}{Decisional Diffie-Hellman (DDH)}{ddh}

  The DDH assumption holds in a group of prime order $p$ with generator $\gen$
if $(\gen^x, \gen^y, \gen^{xy})$ and $(\gen^x, \gen^y, \gen^z)$, for $x, y, z
\sample \Zp$, are computationally indistinguishable.

\end{assumptionbox}

The symmetric external Diffie-Hellman (SXDH) assumption holds in a bilinear group
if DDH holds in both $\group{1}$ and $\group{2}$.

\begin{outline}
  \item DDH cannot hold in a symmetric bilinear group ($\group{1} = \group{2}$), since the pairing decides it; so SXDH needs type III
  \item some schemes we use are only proven in the generic group model (GGM), cite EC:Shoup97, IMA:Maurer05
  \item GGM: adversary sees only opaque handles to group elements, computes only through oracles for the group operations and the pairing
  \item a GGM proof rules out attacks that ignore the representation of the group; weaker than a reduction to a stated assumption
  \item $q$-type assumptions: statement grows with a parameter $q$ (typically the number of signing queries); stronger than fixed-size ones like DDH
\end{outline}

\subsection{Sigma protocols and the Fiat-Shamir transform}\label{sec:prelims-sigma}

\begin{outline}
  \item most proof systems here are Sigma protocols made non-interactive
  \item Sigma protocol for $\rel$: three moves, prover sends $a$, verifier sends uniform challenge $c$, prover sends response $z$, verifier checks against $(\inst, a, c)$
  \item completeness: as for proof systems
  \item special soundness: efficient extractor gets a witness from two accepting transcripts $(a, c, z)$, $(a, c', z')$ with $c \neq c'$
  \item special honest-verifier zero knowledge: efficient simulator, given $\inst$ and $c$, outputs $(a, z)$ distributed as an honest run with challenge $c$ (cite C:CraDamSch94)
\end{outline}

\begin{outline}
  \item example: Schnorr's protocol for knowledge of $x$ with $h = \gen^x$ (cite JC:Schnorr91)
  \item prover sends $a = \gen^r$ for $r \sample \Zp$, gets $c \sample \Zp$, returns $z = r + cx$; verifier checks $\gen^z = a h^c$
  \item same protocol proves a representation $h = \prod_i \gen_i^{x_i}$, one response per exponent; non-interactive form = Schnorr proof
  \item several equations over the same witness: run in parallel with one shared challenge, so they hold together (e.g.\ two group elements with the same discrete logarithm)
\end{outline}

\begin{outline}
  \item Fiat-Shamir: prover computes $c = \hash(\inst, a, \lbl)$ itself, $\hash$ a random oracle, $\lbl$ a label (cite C:FiaSha86)
  \item quasi-unique responses: for given $a$ and $c$, no efficient prover finds two accepting responses; unique responses (stronger): response fixed by a verification equation, as in Schnorr
  \item result: whole statement hashed + special soundness + SHVZK + quasi-unique responses gives a simulation extractable labelled proof system in the ROM (cite INDOCRYPT:FKMV12)
  \item why hash the whole statement: a part left out can be chosen after seeing the challenge
  \item ZK simulator programs the random oracle at $(\inst, a, \lbl)$
  \item extractor rewinds: reruns $\adv$ from the query that fixed the challenge, answers with a fresh challenge, applies special soundness; forking lemma bounds success (cite JC:PoiSte00)
  \item extraction from a single run without rewinding = straight-line/online, what UC needs (cite C:Fischlin05; links back to \Cref{sec:prelims-security})
\end{outline}

\subsection{The Pedersen commitment scheme}

\begin{outline}
  \item \Cref{prim:pedersen}: commits to several field elements at once with one group element (cite C:Pedersen91)
\end{outline}

\begin{primitivebox}{The Pedersen commitment scheme}{pedersen}

  \begin{description}

    \item[$\ck \sample \ComGen(\secparam, n)$ :] On input the security parameter
and message vector length $n$, sample $\gen_1,\ldots,\gen_n, \hgen \sample
\group{1}$ and return $\ck = (\gen_1,\ldots,\gen_n,\hgen)$.

    \item[$\com \leftarrow \ComCommit(\ck,(m_1,\ldots,m_n),\openinfo)$ :] On
input commitment key $\ck$, messages $(m_1,\ldots,m_n) \in \Zp^n$, and opening
$\openinfo \in \Zp$, return $\com = \prod_{i=1}^{n} \gen_i^{m_i}\cdot
\hgen^{\openinfo}$.

    \item[$(m_1,\ldots,m_n) \leftarrow \ComOpen(\ck,\com,\openinfo)$ :] On
input $\ck$, commitment $\com$, and opening $\openinfo = (m_1,\ldots,m_n,
\openinfo')$, return $(m_1,\ldots,m_n)$ if $\com = \prod_{i=1}^{n}
\gen_i^{m_i}\cdot\hgen^{\openinfo'}$, and $\bot$ otherwise.

  \end{description}

\end{primitivebox}

The Pedersen commitment scheme is perfectly hiding, and computationally binding
under the DL assumption~\cite{C:Pedersen91}.

\subsection{The ElGamal encryption scheme}

\begin{outline}
  \item \Cref{prim:elgamal}: encrypts group elements (cite C:ElGamal84)
\end{outline}

\begin{primitivebox}{The ElGamal encryption scheme}{elgamal}

  \begin{description}

    \item[$(\pk, \sk) \sample \PKEgen(\secparam)$ :] Sample $\mu \sample \Zp$,
and return $\sk = \mu$ and $\pk = \gen^{\mu}$.

    \item[$\ciph \sample \PKEencrypt(\pk, M)$ :] On input $M \in \group{1}$,
sample $\rho \sample \Zp$, and return $\ciph = (\gen^{\rho}, M \pk^{\rho})$.

    \item[$M \leftarrow \PKEdecrypt(\sk, \ciph)$ :] On input $\ciph = (f_1, f_2)$,
return $f_2 / f_1^{\mu}$.

  \end{description}

\end{primitivebox}

\begin{outline}
  \item ElGamal is IND-CPA under DDH in $\group{1}$ (cite C:ElGamal84)
  \item messages are group elements: a field element $m$ is encrypted as $\gen^m$, and decryption gives back $\gen^m$, not $m$
\end{outline}

\subsection{Structure-preserving signature schemes}\label{sec:prelims-sps}

\begin{outline}
  \item SPS schemes: signatures where messages, public keys and signatures are vectors of group elements and verification only checks pairing product equations (cite C:AFGHO10, C:AGHO11)
  \item why we want them: nothing is hashed, so they combine with efficient zero-knowledge proofs over the same groups; users prove knowledge of signatures on committed or hidden group elements
\end{outline}

An SPS scheme $\SPS$ over the bilinear groups above consists of the following
algorithms:

\begin{description}

  \item[$(\pk, \sk) \sample \SPSgen(\secparam, \ell)$ :] The key generation
algorithm, on input the security parameter and a vector length $\ell$, outputs
public verification key $\pk$, which consists of group elements, and secret
signing key $\sk$.

  \item[$\SPSsig \sample \SPSsign(\sk, \vec{\msg})$ :] The signing algorithm
signs message $\vec{\msg} \in \group{1}^\ell$ using secret key $\sk$, and outputs
a signature $\SPSsig$ consisting of elements of $\group{1}$ and $\group{2}$.

  \item[$\{0,1\} \leftarrow \SPSverify(\pk, \vec{\msg}, \SPSsig)$ :]
The verification algorithm evaluates pairing product equations in $\pk$,
$\vec{\msg}$ and $\SPSsig$, and outputs $0$ or $1$.

\end{description}
%
Some schemes are rerandomisable: they additionally provide an algorithm
$\SPSsig' \sample \SPSrandomise(\SPSsig)$, which, without the secret key,
turns a signature into a new signature on the same message, distributed as a
freshly generated one.

\begin{outline}
  \item a rerandomisable scheme cannot be strongly unforgeable: rerandomising is itself a strong forgery
\end{outline}

Correctness and EUF-CMA security are defined as for digital signature schemes,
with messages restricted to $\group{1}^\ell$.

\subsection{The $\AGHO$ SPS scheme}

\begin{outline}
  \item Abe, Groth, Haralambiev and Ohkubo: SPS signatures in type III groups have at least three group elements; they give schemes meeting this (cite C:AGHO11)
  \item their main scheme: strongly unforgeable in the GGM; a six-element variant: strongly unforgeable under a non-interactive $q$-type assumption
  \item we use their rerandomisable three-element scheme, $\group{1}$ and $\group{2}$ swapped so messages are in $\group{1}$, and call it $\AGHO$
  \item secret key is only field elements; the message enters signing only through the terms $\msg_i^{-\mu_i}$ (matters for the threshold version in \Cref{chap:tsps})
  \item point to \Cref{prim:agho}
\end{outline}

\begin{primitivebox}{The $\AGHO$ SPS scheme}{agho}

  \begin{description}

    \item[$(\pk, \sk) \sample \SPSgen(\secparam, \ell)$ :] Sample $\mu_1,
\ldots, \mu_\ell, \tau \sample \Zp^*$, and output $\sk = (\mu_1, \ldots, \mu_\ell,
\tau)$ and $\pk = (\widetilde{u}_1, \ldots, \widetilde{u}_\ell, v)$, where
$\widetilde{u}_i = \ggen^{\mu_i}$ and $v = \gen^{\tau}$.

    \item[$\SPSsig \sample \SPSsign(\sk, \vec{\msg})$ :] On input $\vec{\msg} =
(\msg_1, \ldots, \msg_\ell) \in \group{1}^\ell$, sample $\omega \sample \Zp^*$
and output
      %
      \begin{equation*}
      %
	\SPSsig = (\widetilde{r}, \widetilde{s}, z) = \left(\ggen^{\omega}, \
\widetilde{r}^{\,\tau}, \ \left(\gen \prod_{i=1}^{\ell}
\msg_i^{-\mu_i}\right)^{1/\omega}\right).
      %
      \end{equation*}

    \item[$\SPSsig' \sample \SPSrandomise(\SPSsig)$ :] On input $\SPSsig =
(\widetilde{r}, \widetilde{s}, z)$, sample $\omega' \sample \Zp^*$ and output
$(\widetilde{r}^{\,\omega'}, \widetilde{s}^{\,\omega'}, z^{1/\omega'})$.

    \item[$\{0,1\} \leftarrow \SPSverify(\pk, \vec{\msg}, \SPSsig)$ :] On input
$\SPSsig = (\widetilde{r}, \widetilde{s}, z)$, output $1$ if and only if
      %
      \begin{equation*}
      %
	\pairing(v, \widetilde{r}) = \pairing(\gen, \widetilde{s}) \ \wedge\
\pairing(z, \widetilde{r}) \prod_{i=1}^{\ell} \pairing(\msg_i, \widetilde{u}_i)
= \pairing(\gen, \ggen).
      %
      \end{equation*}

  \end{description}

\end{primitivebox}

\begin{outline}
  \item $\AGHO$ is EUF-CMA secure in the GGM (cite C:AGHO11)
  \item being rerandomisable, it is not strongly unforgeable
\end{outline}

\subsection{The BBS+ signature scheme}\label{sec:prelims-bbs}

The BBS+ signature scheme of \Cref{prim:bbs} signs several messages at
once~\cite{C:BonBoySha04,SCN:AuSusMu06,EPRINT:CamDriLeh16}.

\begin{primitivebox}{The BBS+ signature scheme}{bbs}

  \begin{description}

    \item[$(\sk,\pk) \sample \BBSgen(\secparam, n, (\gen_0,\ldots,\gen_{n+1},
\ggen))$ :] On input the security parameter, vector length $n$, and generators
$\gen_0,\gen_1, \ldots, \gen_{n+1} \in \group{1}$, $\ggen \in \group{2}$, do the
following:

      \begin{enumerate}

	\item Sample $z \sample \Zp$.

	\item Return $\sk = z$, $\pk = (\gen_0, \ldots, \gen_{n+1}, \ggen,
\ggen_Z = \ggen^z)$.

      \end{enumerate}

      The generators are an input rather than an output. Standard presentations
have the signer sample them, which is sound when the signer is the only party
relying on them. Where the same generators also key a commitment whose binding
must hold against the signer, they are instead fixed by hashing to $\group{1}$
from a public seed, so that no party knows a discrete logarithm relation among
them.

    \item[$\signature \sample \BBSsign(\sk,(\msg_1, \ldots, \msg_n))$ :] On
input secret key $\sk = z$ and message $(\msg_1, \ldots, \msg_n) \in
\group{1}^n$, do:

      \begin{enumerate}

	\item Sample $e,r \sample \Zp$.

	\item Return $\signature = (e, r, E)$ with $E = (\gen_0 \cdot
\prod_{i=1}^{n} \msg_i \cdot \gen_{n+1}^{r})^{1/(z+e)}$.

      \end{enumerate}

    \item[$\{0,1\} \leftarrow \BBSverify(\pk,(\msg_1, \ldots, \msg_n),
\signature)$ :] On input public key $\pk$, message $(\msg_1, \ldots, \msg_n)$,
and signature $\signature = (e,r,E)$, do:

      \begin{enumerate}

	\item Return $\pairing(E,\ggen_Z\cdot\ggen^{e}) \checkcheck
\pairing(\gen_0 \cdot \prod_{i=1}^{n} \msg_i \cdot \gen_{n+1}^{r}, \ggen)$.

      \end{enumerate}

  \end{description}

\end{primitivebox}

The BBS+ signature scheme is existentially unforgeable under the $q$-strong
Diffie-Hellman ($q$-SDH) assumption of \Cref{asm:qsdh-bbs}, a $q$-type
assumption~\cite{EC:BonBoy04a}.

\begin{assumptionbox}{$q$-Strong Diffie-Hellman ($q$-SDH)}{qsdh-bbs}

  The $q$-SDH assumption holds in the type III bilinear group
$(\group{1},\group{2},\group{T},\pairing)$ if, given $(\gen_0, \gen_0^{\theta},
\ldots, \gen_0^{\theta^{q}}, \ggen, \ggen^{\theta})$ for $\gen_0 \in \group{1}$,
$\ggen \in \group{2}$ and uniformly random $\theta \sample \Zp$, no PPT
adversary can output a pair $\big(\varsigma,\,
\gen_0^{1/(\theta+\varsigma)}\big)$ for any $\varsigma \in \Zp \setminus
\{-\theta\}$ with non-negligible probability.

\end{assumptionbox}

This presentation signs group elements. To sign field elements $m_1, \ldots, m_n
\in \Zp$, the signer signs $M_i = \gen_i^{m_i}$. Signing the group elements
$M_i$ directly, rather than the field elements $m_i$ themselves, is useful
because computing $\prod_i M_i$ requires no knowledge of the discrete logarithms $m_i$,
so a signer can produce a valid signature on committed values without ever
learning them~\cite{SCN:AuSusMu06,EC:TesZhu23b}.

Unforgeability when signing group elements carries a condition that signing
field elements does not. An adversary able to obtain a signature on an arbitrary group
element could ask
for one on $M_i' = M_i \cdot \gen_i^{\delta}$, shifting the certified message by
a quantity it knows. The reduction to $q$-SDH therefore requires the signer to
sign only elements that come with a proof of knowledge of their representation
in the bases $\gen_i$~\cite{SCN:AuSusMu06,EC:TesZhu23b}.

